\section{Сравнение с альтернативными подходами}

\begin{center}
\begin{tabular}{|l|c|c|c|c|}
\hline
Подход & Контекст & Обобщение & Обфускация & Интерпретируемость \\
\hline
Сигнатурный IDS & Нет & Низкое & Низкая & Высокая \\
ML-аномалии & Частично & Среднее & Средняя & Низкая \\
Графы знаний & Да & Среднее & Средняя & Средняя \\
\textbf{Атрибутивные метаграфы} & \textbf{Полный} & \textbf{Высокое} & \textbf{Высокая} & \textbf{Высокая} \\
\hline
\end{tabular}
\end{center}

Предложенный метод сочетает \textbf{интерпретируемость} и \textbf{семантическую глубину}, что делает его пригодным для использования в регулируемых секторах (госсектор, финансы, критическая инфраструктура).